\qhead{7. Costs and Their Relation with Pricing; Financial Implications of Price Changes — What changes in costs have greater implications in terms of price changes for your company? Explain why.}

For pricing only \textbf{incremental} costs matter; sunk R\&D and overhead are irrelevant \parencite{nagle2018}. The cost change with the greatest pricing implication for BYD is the \textbf{battery cell cost}, since the battery is the largest item in an EV's variable bill of materials. The Blade Battery's roughly \euro10/kWh saving \parencite{motleyfool2025} flows straight into the contribution margin and lowers the price floor.

\textbf{Why it dominates — break-even sales change.} A price change is judged by \(\%\,\text{BE} = -\Delta P/(CM+\Delta P)\). Take a model at RMB~100{,}000 with variable cost 80{,}000, so \(CM=20{,}000\) (20\%). A \(-10\%\) cut gives \(\%\,\text{BE}=10{,}000/10{,}000=+100\%\): BYD must \emph{double} volume just to hold profit. But a lower battery cost raising \(CM\) to 30{,}000 cuts that to \(+50\%\). \textbf{Lower variable cost shrinks the break-even volume}, so a cut that is suicidal for a higher-cost rival is merely painful for BYD.

\textbf{Profit leverage.} A 1\% price move typically shifts operating profit 8--11\%. BYD's 2025 net profit fell 19\% on \emph{record} revenue \parencite{byd2025} — thin per-unit concessions across millions of units working the leverage in reverse, partly offset because higher-margin exports cross-subsidise the home-market fight.
